% 5-1-summary.tex
%  Budget: ~2pp.  HANDOFF 16.5, 6.1-6.5, 7.
%  MANDATE:
%   - restate the aim and the significance; introduce NO number that
%     chapters 3-4 did not already report
%   - the nulls and the ties are results, and are summarised as results
%  Rules: decimals in \lr{}; integers Persian; every table cell in \lr{}.

\section{جمع‌بندی نتایج}
\label{sec:5-1}

هدف این پژوهش، پیش‌بینی قیمت روزِ بعدِ برق در بازار آلمان بود، بر دو افقِ ساعتی و
میانگین روزانه، با پنج خانواده‌ی مدل و دو ترکیب، در چارچوبی که بتوان نتیجه‌اش را
هم بازتولید کرد و هم در برابر ادبیات منتشرشده سنجید. \cref{sec:4-7} پاسخ‌ها را
یک‌به‌یک در برابر پرسش‌های پژوهش گذاشت؛ آنچه در ادامه می‌آید، جمع‌بندیِ همان
پاسخ‌هاست و هیچ عدد تازه‌ای به آن‌ها نمی‌افزاید.

\subsection*{آنچه ساخته شد}

خط لوله‌ای کامل از داده‌ی خام تا پیش‌بینیِ امتیازدهی‌شده: بارگذارِ داده با
دریافتِ یک‌باره و ثبتِ اثر انگشت \lr{sha256}، ماتریس ویژگیِ ۲۴۷ ستونی که هیچ
ستونش از پس از مبدأ پیش‌بینی تغذیه نمی‌شود، تنظیم ابرپارامتر با ۵۰ آزمایش بر
پنجره‌ای که کاملاً پیش از دوره‌ی آزمون قرار دارد، و ارزیابیِ پیش‌رونده بر ۷۲۸
مبدأ از ۴ ژانویه ۲۰۱۶ تا ۳۱ دسامبر ۲۰۱۷. سه شیء مدلیِ متمایز \lr{—} پیش‌رونده،
تثبیت‌شده، و برازشِ تفسیری \lr{—} از یکدیگر تفکیک شدند تا هیچ عددی از مسیرِ
نادرست وارد جدول‌ها نشود.

\subsection*{دقت}

بهترین پیکربندی، ترکیب رژیم‌محور بود با \lr{MAE} برابر \lr{3.5569} یورو بر
مگاوات‌ساعت در افق ساعتی، در برابر \lr{7.7501} برای مدل ساده. هر پنج مدل
آموزش‌دیده از مدل ساده بهتر بودند و این برتری پس از تصحیح چندگانگی نیز پابرجا
ماند. اما میان خودِ مدل‌های خوب، فاصله‌ها بسیار کمتر از چیزی است که جدول‌های
مرتب‌شده القا می‌کنند: \lr{LSTM} و \lr{LEAR-LASSO} از یکدیگر تفکیک‌پذیر نیستند
(\lr{p} برابر \lr{0.4036})، و برتریِ عددیِ \lr{LSTM} بر \lr{LightGBM} تصحیح برای
چندگانگی را تاب نمی‌آورد. یعنی در این پیکربندی، هزینه‌ی یادگیری عمیق مزیتِ
قابل‌دفاعی بر یک مدل خطیِ تنظیم‌شده نخرید.

\subsection*{ترکیب و رژیم}

هر دو ترکیب از تک‌تک اعضای خود بهتر بودند، با فاصله‌ای که تصحیح چندگانگی را به
سادگی تاب می‌آورد. وزن‌دهیِ رژیم‌محور نیز ترکیب ایستا را بهبود داد، اما دقیقاً
آنجا که برای آن طراحی شده بود: تمام بهبود بر ۷۷ روز پرتنش متمرکز است
(\lr{p} برابر \lr{0.0063})، و بر ۶۵۱ روز آرام اثر عملاً صفر است
(\lr{p} برابر \lr{0.8465}). نتیجه‌ی تجمیعی \lr{—} \lr{p} برابر \lr{0.0226}
\lr{—} به‌عنوان یک فرضیه‌ی پیش‌اعلام‌شده‌ی تأییدی دفاع می‌شود، و
\cref{sec:4-5-1} صریحاً نوشت که همین عدد اگر از خانواده‌ی اکتشافی برداشته
می‌شد، از تصحیح هولم عبور نمی‌کرد.

\subsection*{افق روزانه}

برای سه مدل از چهار مدل، میانگین‌گیری از پیش‌بینی‌های ساعتی بهتر از آموزشِ
مستقیم بر سری روزانه بود. \lr{SARIMAX} استثناست و در \cref{sec:4-4} استثنا
بودنش توضیح داده شد، نه میانگین گرفته شد.

\subsection*{تفسیرپذیری}

تحلیل \lr{SHAP} نشان داد رفتار مدل میان دو رژیم واقعاً تغییر می‌کند: سهم قیمتِ
روز پیش در شرایط پرتنش به‌شدت بالا می‌رود و سهم پیش‌بینی بار پایین می‌آید. این
یافته از مسیری کاملاً مستقل از آزمون‌های آماری به همان نتیجه‌ای می‌رسد که
وزن‌دهی رژیم‌محور بر آن استوار است، و همین هم‌گرایی است که آن را از یک همبستگیِ
اتفاقی جدا می‌کند.

\subsection*{آنچه به دست نیامد}

سه نتیجه‌ی منفی نیز به همان اندازه بخشی از جمع‌بندی‌اند. نخست، این پژوهش بهترین
مدل منتشرشده‌ی مرجع را شکست نمی‌دهد؛ با دو گونه‌ی آن هم‌تراز است و به یکی
می‌بازد. دوم، برتری یادگیری عمیق بر خط‌مبناهای خطیِ قوی در این پیکربندی
اثبات‌پذیر نبود. سوم، و مهم‌تر از هر دو، مدل‌های تثبیت‌شده در انتقال به داده‌ی
خارج از توزیع شکست خوردند. آن نتیجه ذیل هیچ‌یک از پرسش‌های پژوهش نمی‌گنجد و
موضوع \cref{sec:5-2} است.
